\paragraph{Why did two bus-32 declarations arrive one report late?}
The two late records shared the same six-report pulse and noise
realization, with different governors. Both prefixes already exceeded
the single-terminal reader's threshold against quiet operation. Its best
source alternated between buses 31 and 32. At onset plus 45 reports,
the nominal-governor record favored bus 31 with conditional probability
0.644. In the other record, bus 31 led with probability 0.999 at
plus 44, followed by bus 32 with probability 0.999 at plus 45.
The next scheduled publication was at plus 46. A faster decision clock
alone could therefore have published the wrong bus.

For the revised first-pulse likelihood we kept the three-report clock
and included all seven available current angles alongside the 16
voltage coordinates. Let $H_w$ map rotor-angle changes to these
23 standardized coordinates. Use the original voltage scales and
the same working angular standard deviation of 1~mrad at each
energized current terminal. If $U$ contains the ten orthonormal
columns spanning $H_w$, every candidate response lies in that space
at each report. Projecting each report and template onto $U$ thus
preserves both the inner product and the template energy in
\eqref{eq:pulse-common-prefix}. The residual check still uses all 23
coordinates, including deviations that rotor motion cannot explain.

At the final source comparison we integrate the positive torque
severity $a$ uniformly over $[l,u]=[0.7,1.3]$. For a candidate path
write $p=\langle P_N s,P_N y\rangle$, $e=\|P_N s\|^2>0$, and
$z=p/\sqrt e$. Its marginal likelihood ratio is
\[
 \frac{\sqrt{2\pi/e}}{u-l}\,\exp(z^2/2)
 \big[\Phi(u\sqrt e-z)-\Phi(l\sqrt e-z)\big].
\]
The fast screen retains its three amplitude nodes. In development,
using all currents with three nodes also in the final comparison
produced premature bus-34 declarations on bus-33 pulses. Continuous
severity integration removed these source errors on the same
records. Both late bus-32 cases in the diagnostic replay were
correctly located at plus seven reports, one report after torque
cleared. We excluded the diagnostic replay from the new blind counts.

During recovery, the bank can also announce a new shunt fault while
voltage follows a smooth trajectory. We check this particular
declaration at its proposed publication report. For a candidate
bus $b$, let $\Delta y_b(a)$ be the proposed shunt conductance and
$Z$ the inverse nominal nodal admittance including the generator
Norton admittances. Holding the internal source phasors fixed at
switching gives the rank-one voltage change
\[
 v(a)-v_0=-\frac{\Delta y_b(a)\,v_{0b}}
                    {1+\Delta y_b(a)\,Z_{bb}}\, Z e_b .
\]
The log-magnitudes and phases of $v(a)/v_0$ at the eight PMUs
form the signed jump $j_b(a)$.
At each available onset within two reports of the estimated onset,
subtract the locally extrapolated voltage from the mean of up to
three post-onset reports to give $d$. With $W$ the inverse diagonal
scale estimated from preceding increments, the profiled log likelihood
ratio for a jump against smooth continuation is
\[
 G_b=\tfrac12\left(\|Wd\|^2-\min_a\|W(d-j_b(a))\|^2\right).
\]
The search spans the weak-to-nominal fault range on a grid and
includes its endpoints. We suppress the declaration only when
$G_b\leq -5$ and the smallest admissible jump has
$\|Wj_b(a_{\min})\|\geq5$ at every available onset. Both cutoffs
are operating choices. The check applies only to a fault proposed
by the recovery observer at the current report. In the two
development cohorts it removed two unmatched declarations in one
and none in the other, and retained every previously correct
family--asset declaration in both.

\begin{table}[t]
\caption{Frozen replication on 420 physical events. Det.\ counts on-time triggers,
Fam.\ correct families, and Exact the correct family and asset within 45 reports.
Active counts exact decisions before clearance. Extra counts unmatched declarations
across all 465 streams.}
\label{tab:network-torque-replication}
\centering\footnotesize
\setlength{\tabcolsep}{2.4pt}
\begin{tabular}{lrrrrrr}
\toprule
Method & Macro-F1 & Det. & Fam. & Exact & Active & Extra \\
\midrule
Recovery bank & 0.9853 & 412 & 412 & 405 & 388 & 5 \\
Refitted ExtraTrees & 0.7096 & 302 & 270 & 258 & 116 & 5 \\
One current angle & 0.9896 & 416 & 416 & 408 & 388 & 5 \\
Revised estimator & 0.9896 & 416 & 416 & 409 & 394 & 5 \\
\bottomrule
\end{tabular}
\end{table}


Table~\ref{tab:network-torque-replication} uses 15 new scenario and
noise blocks with matched governor variants. The asset split, eight
PMUs, 30-report/s sampling and 45-report publication deadline match
the preceding experiment. Both estimators located all ten short
bus-32 pulses. The revised estimates arrived four to seven reports
after onset (median six), compared with five to seven for the
single-current estimator (median six). At bus 35 the latter assigned
one pulse to bus 36. The revision located it at report 27 after onset.
Across all 420 events the paired gain in exact diagnoses was 0.24
percentage points (95\% whole-block bootstrap interval 0.00--0.71).
The corresponding gain before clearance was 1.43 points
(0.00--2.86). The unit resampled was the scenario block.

The revised estimator made 120/120 exact fault diagnoses, 112/120
line diagnoses, 117/120 generator diagnoses and 60/60 load diagnoses.
It missed four line events and assigned four other lines and three
generators to the wrong asset. Both estimators made the same five
unmatched declarations. Three occurred after the second line
reclosed. In two governor variants the bank first attributed a
trip of line 10--11 to line 6--11 and then repeated the same
declaration while the line was still open. The shunt check removed
no alarm on these blind records. Of the two spurious fault
declarations, the voltage change at the proposed onset favored a
shunt over local continuation for bus 5. At bus 36 the weakest
shunt had a standardized norm below 1.2 at every tested onset,
too small to rule out. All four methods remained silent
in the 165 quiet streams (1.05 monitored quiet hours).
